% فصل دوم: مبانی نظری و پیشینه پژوهش
\chapter{مبانی نظری و پیشینه پژوهش}
\label{ch:background}

هدف این فصل، مرور مفاهیمی است که برای فهم چارچوب شبکه‌های عصبی آگاه از
فیزیک لازم‌اند. ابتدا معادلات دیفرانسیل جزئی و چند مثال مهم آن‌ها معرفی
می‌شوند و سپس روش‌های عددی کلاسیک حل این معادلات به‌اختصار بررسی می‌گردد.
در ادامه، مبانی شبکه‌های عصبی و آموزش آن‌ها مرور می‌شود و پس از آن به
مشتق‌گیری خودکار، به‌عنوان ابزار کلیدی این پایان‌نامه، می‌پردازیم. در پایان
فصل نیز پیشینه پژوهش در زمینه ترکیب یادگیری ماشین با دانش فیزیک مرور
می‌شود تا جایگاه مقاله مرجع این پایان‌نامه مشخص گردد.

\section{معادلات دیفرانسیل جزئی}
\label{sec:pde}

معادله دیفرانسیل جزئی رابطه‌ای است میان یک تابع مجهول چندمتغیره و مشتق‌های
جزئی آن. به‌طور کلی، یک معادله دیفرانسیل جزئی برای تابع مجهول
$u(\mathbf{x})$ را می‌توان به شکل زیر نوشت:
\begin{equation}
F\left(\mathbf{x}, u, Du, D^2u, \dots, D^ku\right) = 0,
\label{eq:pde-general}
\end{equation}
که در آن $\mathbf{x}$ بردار متغیرهای مستقل (مانند مکان و زمان)، $Du$ شامل
مشتق‌های مرتبه اول، $D^2u$ شامل مشتق‌های مرتبه دوم و به همین ترتیب است.
بزرگ‌ترین مرتبه مشتق ظاهرشده در معادله را مرتبه معادله می‌نامند. اگر تابع
$F$ نسبت به $u$ و مشتق‌های آن خطی باشد، معادله خطی و در غیر این صورت
غیرخطی\LTRfootnote{Nonlinear} نامیده می‌شود \cite{evans2010}.

برای آن‌که مسئله به‌خوبی تعریف شود، معمولاً معادله باید همراه با شرط اولیه
و شرط مرزی در نظر گرفته شود. شرط اولیه مقدار جواب (و گاهی مشتق زمانی آن)
را در لحظه شروع مشخص می‌کند و شرط مرزی رفتار جواب را روی مرز دامنه مکانی
تعیین می‌نماید. مسئله‌ای که در آن جواب با این شرایط اضافی یکتا و نسبت به
داده‌ها پایدار باشد، مسئله خوش‌وضع\LTRfootnote{Well-posed Problem} نامیده
می‌شود.

در ادامه سه مثال کلاسیک را که در این پایان‌نامه نیز به آن‌ها اشاره خواهد
شد معرفی می‌کنیم.

\textbf{معادله گرما.} معادله گرما ساده‌ترین مثال از معادلات سهموی است و
توزیع دما در یک جسم را بر حسب زمان توصیف می‌کند. شکل یک‌بعدی آن چنین است:
\begin{equation}
u_t = \alpha u_{xx},
\label{eq:heat}
\end{equation}
که در آن $u(t,x)$ دما، $\alpha$ ضریب پخش و زیروندها نشان‌دهنده مشتق جزئی
نسبت به زمان ($t$) و مکان ($x$) هستند.

\textbf{معادله برگرز.} معادله برگرز\LTRfootnote{Burgers Equation} یک معادله
غیرخطی است که می‌توان آن را ساده‌شده یک‌بعدی معادلات ناویر ـ استوکس دانست
و به‌همین دلیل در دینامیک سیالات و نظریه ترافیک و غیره کاربرد دارد
\cite{hopf1950}. شکل استاندارد آن عبارت است از:
\begin{equation}
u_t + u u_x - \nu u_{xx} = 0,
\label{eq:burgers-intro}
\end{equation}
که در آن $\nu$ ضریب لزجت است. جمله غیرخطی $uu_x$ باعث تند شدن موج و جمله
$u_{xx}$ باعث پخش و هموار شدن آن می‌شود. رقابت این دو جمله، به‌ویژه به‌ازای
$\nu$های کوچک، جبهه‌های بسیار تندی در جواب ایجاد می‌کند که محاسبه دقیق
آن‌ها برای روش‌های عددی چالشبرانگیز است. به همین سبب، معادله برگرز معیاری
استاندارد برای سنجش روش‌های جدید محسوب می‌شود و در این پایان‌نامه نیز
مسئله نمونه بخش عملی است.

\textbf{معادله شرودینگر.} به‌عنوان مثالی از فیزیک کوانتوم می‌توان معادله
شرودینگر غیرخطی را نام برد که در مقاله مرجع نیز یکی از مسائل نمونه است
\cite{raissi2019}. این مثال نشان می‌دهد که چارچوب مورد بحث به یک حوزه خاص
محدود نیست.

در این پایان‌نامه با دو دسته مسئله سروکار داریم: در مسئله مستقیم، معادله و
پارامترهای آن معلوم‌اند و جواب $u$ مجهول است؛ در مسئله معکوس، بخشی از جواب
به‌صورت مشاهده در اختیار است و هدف، تعیین پارامترهای مجهول معادله است. فصل
سوم هر دو دسته را به‌تفصیل بررسی می‌کند.

\section{روش‌های عددی کلاسیک}
\label{sec:classical}

از آنجا که تنها برای دسته محدودی از معادلات دیفرانسیل جزئی می‌توان جواب
تحلیلی بسته یافت، روش‌های عددی ابزار اصلی حل این معادلات در عمل هستند.
ایده مشترک همه این روش‌ها، تبدیل مسئله پیوسته به یک مسئله گسسته با تعداد
متناهی مجهول است که با کامپیوتر قابل حل باشد \cite{leveque2007}. در ادامه
سه خانواده مهم این روش‌ها را مرور می‌کنیم.

\textbf{روش تفاضل‌های محدود.} در این روش، دامنه با یک شبکه منظم از نقاط
پوشانده می‌شود و مشتق‌ها با ترکیب خطی مقادیر تابع در نقاط مجاور تقریب زده
می‌شوند. برای نمونه، مشتق دوم مکانی با فرمول مرکزی مرتبه دو چنین تقریب
زده می‌شود:
\begin{equation}
u_{xx}(x_i) \approx \frac{u_{i+1} - 2u_i + u_{i-1}}{(\Delta x)^2},
\label{eq:fd}
\end{equation}
که در آن $\Delta x$ فاصله نقاط شبکه است. سادگی پیاده‌سازی مزیت اصلی این
روش است، اما در هندسه‌های پیچیده و برای رسیدن به دقت بالا (که نیازمند
شبکه بسیار ریز است) با مشکل مواجه می‌شود.

\textbf{روش المان‌های محدود.} در این روش، دامنه به قطعه‌های کوچکی به نام
المان تقسیم می‌شود و جواب در هر المان با ترکیب چند تابع پایه ساده تقریب
زده می‌شود. سپس با استفاده از فرمول‌بندی ضعیف\LTRfootnote{Weak Formulation}
معادله، یک دستگاه معادلات جبری برای ضرایب مجهول به دست می‌آید. قدرت اصلی
این روش در برخورد با هندسه‌های پیچیده و شرایط مرزی متنوع است و به همین
دلیل در نرم‌افزارهای مهندسی کاربرد گسترده‌ای دارد.

\textbf{روش‌های طیفی.} در روش‌های طیفی، جواب به‌صورت ترکیب خطی از توابع
پایه سراسری مانند چندجمله‌ای‌ها یا توابع مثلثاتی نوشته می‌شود. برای مسائل
با جواب هموار، این روش‌ها همگرایی بسیار سریعی (نمایی) دارند؛ یعنی با تعداد
نقاط نسبتاً کم به دقت بالایی می‌رسند \cite{trefethen2000}. در بخش عملی این
پایان‌نامه نیز از یک روش طیفی بر پایه تبدیل فوریه برای تولید جواب مرجع
معادله برگرز استفاده شده است.

مقایسه اجمالی این سه خانواده روش در جدول \ref{tab:numerical} آمده است.
نکته مشترکی که همه روش‌های کلاسیک دارند، وابستگی به شبکه محاسباتی و نیاز
به اجرای مجدد از ابتدا به‌ازای هر شرط اولیه یا پارامتر جدید است. همچنین در
ابعاد بالا، تعداد نقاط شبکه به‌صورت نمایی رشد می‌کند که همان پدیده نفرین
ابعاد است. این محدودیت‌ها انگیزه اصلی جست‌وجو برای روش‌های جایگزین، از
جمله روش‌های مبتنی بر شبکه‌های عصبی، به شمار می‌رود.

\begin{table}[t]
\caption{مقایسه اجمالی روش‌های عددی کلاسیک حل معادلات دیفرانسیل جزئی}
\label{tab:numerical}
\centering
\singlespacing
\begin{tabular}{|c|p{3.4cm}|p{3.4cm}|p{3.4cm}|}
\hline
\textbf{روش} & \textbf{ایده اصلی} & \textbf{مزیت مهم} & \textbf{محدودیت مهم} \\
\hline
تفاضل‌های محدود & تقریب مشتق با مقادیر نقاط مجاور روی شبکه منظم & سادگی پیاده‌سازی و تحلیل & دشواری در هندسه‌های پیچیده \\
\hline
المان‌های محدود & تقسیم دامنه به المان و استفاده از فرمول‌بندی ضعیف & انعطاف در هندسه و شرط مرزی & پیچیدگی پیاده‌سازی و هزینه محاسباتی \\
\hline
روش‌های طیفی & بسط جواب بر حسب توابع پایه سراسری & دقت بسیار بالا برای جواب‌های هموار & محدودیت در هندسه‌های ساده و جواب‌های ناهموار \\
\hline
\end{tabular}
\end{table}

\section{شبکه‌های عصبی مصنوعی}
\label{sec:ann}

شبکه عصبی مصنوعی\LTRfootnote{Artificial Neural Network (ANN)} مدلی محاسباتی
است که از ساختار مغز الهام گرفته شده و از واحدهای ساده‌ای به نام نورون
تشکیل می‌شود. هر نورون، ترکیب خطی ورودی‌های خود را محاسبه کرده و سپس یک
تابع غیرخطی به نام تابع فعال‌ساز\LTRfootnote{Activation Function} روی آن اعمال
می‌کند:
\begin{equation}
y = \sigma\left(\sum_{i=1}^{n} w_i x_i + b\right),
\label{eq:neuron}
\end{equation}
که در آن $x_i$ ورودی‌ها، $w_i$ وزن‌ها، $b$ بایاس و $\sigma$ تابع فعال‌ساز
است. اتصال لایه‌هایی از این نورون‌ها به یکدیگر، شبکه‌ای به نام پرسپترون
چندلایه\LTRfootnote{Multi-Layer Perceptron (MLP)} می‌سازد. خروجی لایه $l$ام
چنین شبکه‌ای به‌صورت برداری چنین نوشته می‌شود:
\begin{equation}
\mathbf{h}^{(l)} = \sigma\left(\mathbf{W}^{(l)} \mathbf{h}^{(l-1)} + \mathbf{b}^{(l)}\right),
\label{eq:mlp}
\end{equation}
که $\mathbf{h}^{(0)}$ همان بردار ورودی است. نمایی شماتیک از یک پرسپترون
چندلایه با دو لایه پنهان در شکل \ref{fig:mlp} نشان داده شده است.

\begin{figure}[t]
\centering
\begin{latin}
\begin{tikzpicture}[>=Stealth, node distance=1.6cm,
  neuron/.style={circle, draw, minimum size=7mm, inner sep=0pt},
  ann/.style={->, thin}]
  % input layer
  \node[neuron] (i1) {$x_1$};
  \node[neuron, below=0.5cm of i1] (i2) {$x_2$};
  % hidden layer 1
  \node[neuron, right=1.8cm of i1, yshift=0.6cm] (h11) {};
  \node[neuron, below=0.5cm of h11] (h12) {};
  \node[neuron, below=0.5cm of h12] (h13) {};
  % hidden layer 2
  \node[neuron, right=1.8cm of h11] (h21) {};
  \node[neuron, below=0.5cm of h21] (h22) {};
  \node[neuron, below=0.5cm of h22] (h23) {};
  % output
  \node[neuron, right=1.8cm of h22] (o1) {$y$};
  % connections
  \foreach \s in {i1,i2} \foreach \d in {h11,h12,h13} \draw[ann] (\s) -- (\d);
  \foreach \s in {h11,h12,h13} \foreach \d in {h21,h22,h23} \draw[ann] (\s) -- (\d);
  \foreach \s in {h21,h22,h23} \draw[ann] (\s) -- (o1);
  % labels
  \node[above=0.2cm of h12, font=\small] {Hidden layers};
  \node[above=0.2cm of i1, font=\small] {Input};
  \node[above=0.2cm of o1, font=\small] {Output};
\end{tikzpicture}
\end{latin}
\caption{نمایی شماتیک از یک پرسپترون چندلایه با دو ورودی، دو لایه پنهان و یک خروجی}
\label{fig:mlp}
\end{figure}

انتخاب تابع فعال‌ساز اثر مهمی بر رفتار شبکه دارد. پرکاربردترین توابع
فعال‌ساز در جدول \ref{tab:activations} آمده‌اند. در شبکه‌های آگاه از فیزیک
معمولاً از تابع تانژانت هیپربولیک استفاده می‌شود، زیرا هموار و بی‌نهایت‌بار
مشتق‌پذیر است و مشتق‌گیری‌های پی‌درپی نسبت به ورودی‌ها (که در فصل سوم خواهیم
دید ضروری است) برای آن به‌خوبی تعریف می‌شود؛ در حالی که توابعی مانند
\lr{ReLU} در مبدأ مشتق ندارند و مشتق دوم آن‌ها تقریباً همه‌جا صفر است.

\begin{table}[t]
\caption{توابع فعال‌ساز پرکاربرد در شبکه‌های عصبی}
\label{tab:activations}
\centering
\singlespacing
\begin{tabular}{|c|c|c|}
\hline
\textbf{نام تابع} & \textbf{فرمول} & \textbf{ویژگی} \\
\hline
سیگموئید & $\sigma(z) = \dfrac{1}{1+e^{-z}}$ & خروجی در بازه $(0,1)$، هموار \\
\hline
تانژانت هیپربولیک & $\tanh(z) = \dfrac{e^{z}-e^{-z}}{e^{z}+e^{-z}}$ & خروجی در بازه $(-1,1)$، هموار \\
\hline
واحد خطی یکسوشده & $\mathrm{ReLU}(z) = \max(0,z)$ & محاسبه سریع، نامشتق‌پذیر در صفر \\
\hline
\end{tabular}
\end{table}

توانایی نظری شبکه‌های عصبی در تقریب توابع، با قضیه تقریب جهانی تضمین
می‌شود. صورت ساده‌ای از این قضیه که توسط سیبنکو\LTRfootnote{Cybenko} اثبات شده
\cite{cybenko1989} چنین بیان می‌شود:

\begin{theorem}[تقریب جهانی]
هر تابع پیوسته روی یک مجموعه فشرده را می‌توان با یک شبکه عصبی پیش‌خور با
یک لایه پنهان و تابع فعال‌ساز غیرخطی مناسب، با هر دقت دلخواه تقریب زد.
\end{theorem}

البته این قضیه فقط وجود چنین شبکه‌ای را تضمین می‌کند و درباره تعداد نورون‌های
لازم یا چگونگی یافتن وزن‌ها چیزی نمی‌گوید. در عمل، شبکه‌های عمیق‌تر (با
لایه‌های بیشتر) معمولاً با تعداد پارامتر کمتری به دقت مشابه می‌رسند و به
همین دلیل در کاربردها ترجیح داده می‌شوند.

\section{آموزش شبکه‌های عمیق}
\label{sec:training}

آموزش شبکه عصبی یعنی یافتن وزن‌ها و بایاس‌هایی که خروجی شبکه را به خروجی
مطلوب نزدیک کند. برای این کار ابتدا یک تابع هزینه\LTRfootnote{Loss Function}
تعریف می‌شود که فاصله پیش‌بینی شبکه از مقادیر واقعی را اندازه می‌گیرد.
پرکاربردترین تابع هزینه در مسائل رگرسیون، میانگین مربعات خطا\LTRfootnote{Mean Squared Error (MSE)} است:
\begin{equation}
\mathrm{MSE} = \frac{1}{N}\sum_{i=1}^{N} \left(u_\theta(\mathbf{x}_i) - u_i\right)^2,
\label{eq:mse}
\end{equation}
که در آن $u_\theta$ خروجی شبکه با پارامترهای $\theta$ و $u_i$ مقدار واقعی
نظیر ورودی $\mathbf{x}_i$ است.

کمینه‌سازی تابع هزینه معمولاً با روش‌های مبتنی بر گرادیان انجام می‌شود.
ساده‌ترین آن‌ها، روش گرادیان کاهشی است که در هر گام، پارامترها را در خلاف
جهت گرادیان تابع هزینه به‌روز می‌کند:
\begin{equation}
\theta_{k+1} = \theta_k - \eta \nabla_\theta \mathcal{L}(\theta_k),
\label{eq:gd}
\end{equation}
که $\eta$ نرخ یادگیری\LTRfootnote{Learning Rate} و $\mathcal{L}$ تابع هزینه
است. محاسبه کارای گرادیان در شبکه‌های بزرگ با الگوریتم پس‌انتشار\LTRfootnote{Backpropagation}
انجام می‌شود که در واقع کاربرد قاعده زنجیره‌ای مشتق روی گراف محاسباتی شبکه
است که توسط رومل‌هارت و همکارانش معرفی و فراگیر شد \cite{rumelhart1986}.

دو بهینه‌ساز پرکاربرد که در این پایان‌نامه نیز از آن‌ها استفاده شده عبارت‌اند از:

\textbf{روش آدام.} آدام\LTRfootnote{Adaptive Moment Estimation (Adam)} یکی از
محبوب‌ترین بهینه‌سازها در یادگیری عمیق است که برای هر پارامتر، نرخ یادگیری
تطبیقی جداگانه‌ای نگه می‌دارد \cite{kingma2015}. قاعده به‌روزرسانی آن با
استفاده از میانگین‌های متحرک گرادیان ($m$) و مربع گرادیان ($v$) چنین است:
\begin{equation}
\begin{aligned}
m_k &= \beta_1 m_{k-1} + (1-\beta_1) g_k, \\
v_k &= \beta_2 v_{k-1} + (1-\beta_2) g_k^2, \\
\theta_{k+1} &= \theta_k - \eta \frac{\hat{m}_k}{\sqrt{\hat{v}_k} + \epsilon},
\end{aligned}
\label{eq:adam}
\end{equation}
که در آن $g_k$ گرادیان در گام $k$ام است و $\hat{m}_k$ و $\hat{v}_k$ نسخه‌های
اصلاح بایاس‌شده $m_k$ و $v_k$ هستند. مقادیر پیش‌فرض $\beta_1=0.9$،
$\beta_2=0.999$ و $\epsilon=10^{-8}$ در بیشتر کاربردها به‌خوبی عمل می‌کنند.

\textbf{روش ال-بی‌اف‌جی‌اس.} این روش از خانواده روش‌های شبه‌نیوتنی است که با
تقریب معکوس ماتریس هشین، از اطلاعات انحنای تابع هزینه نیز بهره می‌برد و
معمولاً در نزدیکی نقطه بهینه همگرایی سریعی دارد \cite{liu1989}. نسخه
حافظه‌محدود آن (\lr{L-BFGS}) برای مسائل با پارامترهای زیاد مناسب است. در
آموزش شبکه‌های آگاه از فیزیک، راهبرد رایج این است که ابتدا چند هزار گام با
آدام پیش می‌روند و سپس برای دقت نهایی، آموزش را با \lr{L-BFGS} ادامه
می‌دهند؛ ما نیز در پیاده‌سازی خود همین راهبرد را به کار برده‌ایم.

نکته مهم دیگر، مقداردهی اولیه وزن‌هاست. اگر وزن‌ها نامناسب مقداردهی شوند،
آموزش یا بسیار کند می‌شود یا اصلاً همگرا نمی‌شود. روش گزاویه\LTRfootnote{Xavier Initialization}
که واریانس وزن‌های اولیه را متناسب با تعداد ورودی و خروجی هر لایه تنظیم
می‌کند \cite{glorot2010}، انتخابی استاندارد برای شبکه‌هایی با فعال‌ساز
تانژانت هیپربولیک است و در این پژوهش نیز از آن استفاده شده است.

\section{مشتق‌گیری خودکار}
\label{sec:autodiff}

مشتق‌گیری خودکار مجموعه‌ای از تکنیک‌ها برای محاسبه دقیق مشتق توابعی است که
با یک برنامه کامپیوتری پیاده‌سازی شده‌اند \cite{baydin2018}. تفاوت آن با
مشتق‌گیری عددی (مانند تفاضل محدود) در این است که خطای برشی ندارد و نتیجه
تا دقت ماشین دقیق است؛ و تفاوت آن با مشتق‌گیری نمادی در این است که با
عبارت‌های ریاضی سروکار ندارد، بلکه روی گراف محاسباتی برنامه و با اعمال
قاعده زنجیره‌ای روی عملیات پایه عمل می‌کند، پس با مشکل انفجار عبارت‌ها هم
مواجه نیست.

دو حالت اصلی مشتق‌گیری خودکار عبارت‌اند از حالت مستقیم\LTRfootnote{Forward Mode}
و حالت معکوس\LTRfootnote{Reverse Mode}. در حالت معکوس که همان ایده حاکم بر
الگوریتم پس‌انتشار است، ابتدا برنامه به‌ترتیب اجرا و مقادیر میانی ذخیره
می‌شوند و سپس مشتق‌ها از خروجی به سمت ورودی منتشر می‌گردند. مزیت بزرگ این
حالت آن است که گرادیان یک خروجی اسکالر نسبت به هزاران ورودی (مثلاً وزن‌های
شبکه) تنها با حدود دو برابر هزینه یک اجرای مستقیم به دست می‌آید؛ به همین
دلیل ستون فقرات آموزش شبکه‌های عمیق محسوب می‌شود.

اهمیت مشتق‌گیری خودکار در این پایان‌نامه دو جنبه دارد: اولاً برای به‌روزرسانی
وزن‌های شبکه در فرایند آموزش به آن نیاز داریم (همان کاربرد معمول) و ثانیاً،
و مهم‌تر، برای محاسبه مشتق‌های جواب نسبت به متغیرهای ورودی شبکه یعنی مکان
و زمان. همان‌طور که در فصل سوم خواهیم دید، همین قابلیت است که اجازه می‌دهد
باقیمانده معادله دیفرانسیل را مستقیماً از روی شبکه بسازیم، بدون آن‌که نیازی
به شبکه محاسباتی یا فرمول تفاضل محدود باشد. کتابخانه‌های مدرن یادگیری عمیق
مانند تنسورفلو \cite{abadi2016} و جکس \cite{jax2018} این قابلیت را به‌صورت
توکار فراهم می‌کنند.

\section{پیشینه پژوهش}
\label{sec:related}

ایده استفاده از شبکه‌های عصبی برای حل معادلات دیفرانسیل قدمتی بیش از سه
دهه دارد. لاگاریس و همکارانش در سال ۱۹۹۸ نشان دادند که می‌توان جواب معادلات
دیفرانسیل معمولی و جزئی را با یک شبکه عصبی تقریب زد، به‌گونه‌ای که شرایط
مرزی به‌صورت ساختاری در پاسخ پیشنهادی لحاظ شود و باقیمانده معادله در تعدادی
نقطه نمونه کمینه گردد \cite{lagaris1998}. هرچند این کار در زمان خود مورد
توجه قرار گرفت، اما محدودیت توان محاسباتی و نبود ابزارهای مشتق‌گیری خودکار
کارامد باعث شد این مسیر سال‌ها کم‌رونق بماند.

موج دوم این پژوهش‌ها با ورود فرایندهای گاوسی\LTRfootnote{Gaussian Process} به
این حوزه شکل گرفت. رئیسی و همکارانش نشان دادند که با طراحی کرنل‌های سازگار
با عملگر دیفرانسیل می‌توان جواب معادلات خطی \cite{raissi2017a} و سپس با
خطی‌سازی موضعی، معادلات غیرخطی \cite{raissi2018} را استنتاج کرد؛ روشی که
علاوه بر جواب، برآورد عدم قطعیت هم ارائه می‌داد. اما ماهیت بیزی این روش‌ها
و نیاز به خطی‌سازی، کاربرد آن‌ها را در مسائل به‌شدت غیرخطی محدود می‌کرد.

در مسیری موازی، کارپاتنه و همکارانش چارچوب «علم داده هدایت‌شده با نظریه» را
مطرح کردند که هدف آن سازگار کردن مدل‌های یادگیری ماشین با دانش علمی موجود
در حوزه مسئله است \cite{karpatne2017}. همچنین پژوهش‌های اشمیت و لیپسون
\cite{schmidt2009} و روش \lr{SINDy} معرفی‌شده توسط برانتون و همکارانش
\cite{sindy2016} نشان دادند که می‌توان قوانین حاکم بر یک سیستم دینامیکی را
مستقیماً از روی داده کشف کرد؛ هرچند این روش‌ها معمولاً به کتابخانه‌ای از
جملات کاندید نیاز دارند و برای معادلات جزئی پیچیده چندان مقیاس‌پذیر نیستند.

در این میان، مقاله رئیسی، پردیکاریس و کارنیاداکیس در سال ۲۰۱۹ نقطه عطفی
محسوب می‌شود \cite{raissi2019}. آن‌ها با ترکیب شبکه‌های عصبی عمیق و
مشتق‌گیری خودکار، چارچوبی یکپارچه برای هر دو دسته مسائل مستقیم (یافتن جواب)
و معکوس (کشف پارامترها) ارائه کردند که نیازی به خطی‌سازی، شبکه محاسباتی یا
کتابخانه جملات کاندید ندارد. سادگی ایده، عمومیت آن و نتایج قانع‌کننده روی
مسائل کلاسیک (از معادله برگرز و شرودینگر تا ناویر ـ استوکس) باعث شد این
مقاله به‌سرعت به یکی از آثار مرجع حوزه تبدیل شود و صدها پژوهش بعدی را رقم
بزند؛ به‌گونه‌ای که امروزه مرورهای جامعی مانند \cite{cuomo2022} صرفاً به
دسته‌بندی و تحلیل این خانواده از روش‌ها اختصاص یافته است. این پایان‌نامه
نیز بر مبنای همین مقاله مرجع بنا شده و می‌کوشد ایده‌های آن را در سطح
کارشناسی، هم به‌صورت نظری و هم به‌صورت عملی، بازتولید و ارزیابی کند.

\section{جمع‌بندی فصل}
\label{sec:ch2-summary}

در این فصل، مفاهیم پایه مورد نیاز این پژوهش مرور شد: معادلات دیفرانسیل جزئی
و مثال‌های مهم آن‌ها، روش‌های عددی کلاسیک و محدودیت‌هایشان، ساختار و آموزش
شبکه‌های عصبی، مشتق‌گیری خودکار و پیشینه پژوهش در زمینه یادگیری ماشین آگاه
از فیزیک. در فصل بعد، با تکیه بر همین مبانی، چارچوب شبکه‌های عصبی آگاه از
فیزیک به‌تفصیل تشریح می‌شود.
